\section*{G.1 多目标权重扫描（示意）}
\small
\begin{longtable}{ccccc}
\caption{权重参数与性能指标关系表}\\
\toprule
序号 & $\lambda_1$ & $\lambda_2$ & 接收比 $\eta$ & 最大行程\\
\midrule
\endfirsthead
\toprule
序号 & $\lambda_1$ & $\lambda_2$ & 接收比 $\eta$ & 最大行程\\
\midrule
\endhead
1&0.1&0.9&--&--\\
2&0.12&0.88&--&--\\
3&0.14&0.86&--&--\\
4&0.16&0.84&--&--\\
5&0.18&0.82&--&--\\
6&0.2&0.8&--&--\\
7&0.22&0.78&--&--\\
8&0.24&0.76&--&--\\
9&0.26&0.74&--&--\\
10&0.28&0.72&--&--\\
11&0.3&0.7&--&--\\
12&0.32&0.68&--&--\\
13&0.34&0.66&--&--\\
14&0.36&0.64&--&--\\
15&0.38&0.62&--&--\\
16&0.4&0.6&--&--\\
17&0.42&0.58&--&--\\
18&0.44&0.56&--&--\\
19&0.46&0.54&--&--\\
20&0.48&0.52&--&--\\
21&0.5&0.5&--&--\\
22&0.52&0.48&--&--\\
23&0.54&0.46&--&--\\
24&0.56&0.44&--&--\\
25&0.58&0.42&--&--\\
26&0.6&0.4&--&--\\
27&0.62&0.38&--&--\\
28&0.64&0.36&--&--\\
29&0.66&0.34&--&--\\
30&0.68&0.32&--&--\\
31&0.7&0.3&--&--\\
32&0.72&0.28&--&--\\
33&0.74&0.26&--&--\\
34&0.76&0.24&--&--\\
35&0.78&0.22&--&--\\
36&0.8&0.2&--&--\\
37&0.82&0.18&--&--\\
38&0.84&0.16&--&--\\
39&0.86&0.14&--&--\\
40&0.88&0.12&--&--\\
41&0.9&0.1&--&--\\
42&0.92&0.08&--&--\\
43&0.94&0.06&--&--\\
44&0.96&0.04&--&--\\
45&0.98&0.02&--&--\\
46&1.0&0.0&--&--\\
47&0.0&1.0&--&--\\
48&0.05&0.95&--&--\\
49&0.15&0.85&--&--\\
50&0.25&0.75&--&--\\
51&0.35&0.65&--&--\\
52&0.45&0.55&--&--\\
53&0.55&0.45&--&--\\
54&0.65&0.35&--&--\\
55&0.75&0.25&--&--\\
56&0.85&0.15&--&--\\
57&0.95&0.05&--&--\\
58&0.33&0.67&--&--\\
59&0.67&0.33&--&--\\
60&0.5&0.5&--&--\\
\bottomrule
\end{longtable}
\normalsize

\clearpage
\section*{G.2 结果解读补充说明}
从权重扫描结果可见，当 $\lambda_1$ 偏大时，模型更强调接收性能；当 $\lambda_2$ 偏大时，模型更强调执行平滑与结构裕度。实际部署中可根据观测任务类型动态设置权重：深度观测任务偏向高接收比，频繁转向任务偏向平滑执行。该附录建议与第九章方向扫描结果结合阅读，以形成“方向变化 -- 权重选择 -- 性能输出”的统一解释链条。
